\section{NormCode as a Reasoning Medium}
\label{sec:normcode}

A reasoning medium is a representational language in which reasoning plans are
formulated, not merely described. This section demonstrates that NormCode satisfies six
properties that together constitute a working reasoning medium for AI orchestration.
The formal language specification is in \cite{anon2025}.

\subsection{Six Properties of a Working Reasoning Medium}

\paragraph{P1 --- Interpretable.}
A reasoning medium must be legible to humans at multiple roles without programming
expertise. The \texttt{.ncds} authoring format uses three markers and indentation ---
syntax that reads like structured natural language. The compiler's Formalization phase
produces \texttt{.ncn} (NormCode Natural): a plain-prose narrative available before any
LLM call. A manager, compliance officer, or domain expert reads the \texttt{.ncn} and
verifies the reasoning structure matches intent (C2). Because concept names are
natural-language identifiers, plans can be authored in any language the authoring team
uses: the production PPT Generation plan uses Chinese concept names throughout, and the
scope rule, compiler, and orchestrator operate identically regardless of the natural
language used.

\paragraph{P2 --- Enforceable.}
What is reviewed must be exactly what executes. The compiler determines plan structure
and each step's scope before the first run; the orchestrator enforces the scope rule at
runtime by constructing each step's input from only its declared references. The
\texttt{.ncn} is a deterministic rendering of the actual formal plan --- not a
paraphrase. The medium has structural teeth.

\paragraph{P3 --- Composable.}
Three symbols and one scope rule combine into arbitrarily complex workflows: nested
loops, conditionals, parallel derivation, multi-agent coordination. A 5-node plan and a
40-node plan are governed by the same grammar.

\paragraph{P4 --- Locally Addressable.}
Every reasoning step has a dot-separated flow index (1.1, 1.3.2, etc.). The scope rule
makes every step epistemically self-contained: a developer can examine step 1.3.2
without understanding steps 1.1--1.3.1, because the scope rule guarantees 1.3.2 depends
only on what is declared in its own block.

\paragraph{P5 --- Generalizable.}
The three primitives are domain-agnostic. The same syntax expresses a PPT generation
plan, a code review pipeline, a BOM matching task, a scientific data analysis, or a
medical treatment pathway --- without domain-specific extensions.

\paragraph{P6 --- Portable.}
A NormCode plan is a plain text \texttt{.ncds} file accompanied by
\texttt{concept\_repo.json} and \texttt{inference\_repo.json}. These artifacts are
version-controlled, shareable, and executable by any compliant orchestrator.

\subsection{Syntax and the Scope Rule}

Three markers constitute the full syntax:
\begin{itemize}
  \item \texttt{<-}\ \ Value Concept: A named value, document, or intermediate result
    --- data flowing between steps.
  \item \texttt{<=}\ \ Functional Concept: An operation that derives its parent concept
    from declared inputs.
  \item \texttt{<*}\ \ Context Concept: In-loop state: marks the current element during
    iteration.
\end{itemize}

\textbf{Indentation defines scope.} The inputs available to any inference are exactly
the concepts declared in its immediately enclosing indented block. There is no implicit
sharing. This is the scope rule: the compiler determines each step's scope during
formalization; the orchestrator enforces it at runtime by constructing each step's
input from only its declared references. If an inference requires data not present in
its scope, execution proceeds with the available inputs --- the remedy is to revise the
plan to declare the missing dependency explicitly.

A two-concept example (\texttt{.ncds} authoring format):

\begin{lstlisting}
<- summary
    <= summarize the findings
    <- report
        <= read the uploaded file
        <- source_doc
    <- style_guide
\end{lstlisting}

Read bottom-up: \texttt{source\_doc} is read to produce \texttt{report};
\texttt{report} and \texttt{style\_guide} are the declared inputs to the summarization
inference. The \texttt{style\_guide} at the \texttt{summary} scope is explicit --- not
inherited.

\textbf{Flow indices} assigned during Formalization give each concept a stable address.
Following the sibling pattern --- the functional concept is always \texttt{.1} under its
parent, value inputs are siblings at \texttt{.2}, \texttt{.3}, etc.:
\texttt{summary} $\rightarrow$ 1, \texttt{report} $\rightarrow$ 1.2,
\texttt{source\_doc} $\rightarrow$ 1.2.2, \texttt{style\_guide} $\rightarrow$ 1.3.
Every checkpoint is stored and retrieved by (\texttt{run\_id}, \texttt{flow\_index}).

\textbf{Iteration.} The context concept (\texttt{<*}) marks the current element during
a loop:

\begin{lstlisting}
<- all summaries
    <= for each document in the collection
    <- summary
        <= summarize this document
        <- document
    <- documents
    <* document
\end{lstlisting}

Each iteration of \texttt{summary} receives only its own \texttt{document} instance;
sibling iterations are invisible to each other.

\subsection{Semantic vs.\ Syntactic Separation}

\textbf{Semantic inferences} invoke LLMs to create new information. \textit{Imperative}
inferences produce a concept (a named-axis tensor); \textit{Judgement} inferences
evaluate a truth condition and return a boolean. Both are probabilistic and consume
tokens.

\textbf{Syntactic inferences} perform deterministic data manipulation without LLM
involvement using four operator families: \textit{Assigning} (\$) routes data between
concepts; \textit{Grouping} (\&) collects or concatenates concepts; \textit{Timing}
(@) gates execution on conditions or sequencing dependencies; \textit{Looping} (*)
manages iteration with carried state. These are instant, zero-cost, and fully
deterministic.

In production plans, 60--70\% of nodes are syntactic. The Canvas App renders semantic
nodes as purple hexagons and syntactic nodes as gray rectangles. Syntactic checkpoints
are perfectly reproducible cases (same input $\rightarrow$ same output, always);
semantic checkpoints are bounded-stochastic.

\subsection{The Compilation Pipeline}

A NormCode plan is compiled through four phases:
\begin{enumerate}
  \item \textbf{Derivation} (LLM-assisted): extracts concepts, inference operations,
    dependencies, and hierarchical structure from the \texttt{.ncds} draft.
  \item \textbf{Formalization} (rule-based): assigns flow indices; determines sequence
    types; generates the \texttt{.ncn} plain-prose narrative.
  \item \textbf{Post-Formalization}: assigns paradigm IDs (execution configurations),
    tool faculties, tensor axis declarations, and agent assignments.
  \item \textbf{Activation}: extracts \texttt{concept\_repo.json} and
    \texttt{inference\_repo.json} for the orchestrator.
\end{enumerate}

The \texttt{.ncn} output of Phase~2 is the pre-execution interpretability artifact
(C2). For the slide generation plan:

\begin{lstlisting}
(OUTPUT) collection of written slides
    (ACTION) is obtained by iterating over slide-outlines, where each iteration:
        (ACTION) is returning the following as individual output:
        (VALUE) slide-page-written
            (ACTION) is to write the following:
            (VALUE) slide-page
                (ACTION) is obtained by generating page from
                (VALUE) topic
                (VALUE) style_guide
                (VALUE) current_slide_outline
    (VALUE) slide-outlines
    (CONTEXT) for each current_slide_outline in the iteration
\end{lstlisting}

\subsection{The Reference System and Perceptual Signs}

Concepts' data are stored as \textbf{References} --- $N$-dimensional named-axis tensors
whose entries can be values, file pointers, or \textit{perceptual signs}: compact
symbolic pointers of the form \texttt{\%\{norm\_type\}unique\_id(signifier)} that
identify data without loading content into the prompt. Large inputs pass between steps
as tokens, preventing token-level context accumulation alongside the structural
isolation enforced by the scope rule.

When a step requires actual content, the orchestrator \textit{transmutes} the perceptual
sign at the moment of the LLM call --- never earlier. The Tensor Inspector visualizes
References in table, list, or JSON view --- making every completed step's exact bounded
inputs inspectable in O(1) operations (C1).

\textbf{Environmental replication.} A checkpoint whose tensors contain perceptual signs
is structurally self-contained, but resuming from it requires the referenced external
resources (files, APIs, databases) to be accessible and consistent. This is a
precondition for case reuse, distinct from the structural isolation guarantee: the plan
encodes \textit{what} was used; whether that environment is still available is outside
the language's control. See Sect.~\ref{sec:discussion}.
